\bookchapter{}{Technical preface}{How the book is built, and how it goes with the simulator.}
\label{ch:preface}

\lead{This book is a companion. It was made together with a simulator, Anemostatos, and the two follow each other step by step: the same lessons, with the same numbers and the same titles, the same experiments, the same colours. The simulator lets you watch a feedback loop work, change it and break it; the book explains what you see. Each is incomplete without the other.}

\section*{A companion to the simulator}
Every chapter of the book is one lesson of the application. A lesson there sets up an experiment — a drone, a controller, a wind, a sensor — and starts it; the chapter here sets up the same experiment and asks what will happen and why. Every time series in the book was flown by the simulator itself, with the same code and the same set-up as the lesson in the application, so a plot on the page and a chart on the screen are two views of one flight. Dark boxes styled like the application say where an idea lives on the screen and in the source code; experiments marked \emph{Try this} send you back to the simulator with a question to answer first. The book can be read without the simulator open, but it was written to be read next to it.

\section*{From the particular to the general}
Control theory has a reputation for being abstract. Its textbooks usually begin with the general — Laplace transforms, transfer functions, pole-zero maps — and come to a real machine late, if at all. This book is built the other way round, and that is its main idea: it goes \emph{from the detail and from practice to the general}.

Every chapter starts from one particular thing that can be watched: one drone, one experiment, one number on the screen — a drone that sags below its setpoint, bounces, overshoots, shakes or flips. Then it asks \emph{why}, and works out the answer for this drone, with its real mass and its real gains, building exactly as much mathematics as the answer needs: no more, but also no less. Only then does it step back. The section \emph{The engineer's view} says the same result in general terms, for any system of the same kind; the subsection \emph{In formal terms} states it exactly, as definitions and theorems with their hypotheses and proofs; and the page \emph{Out in the World} shows the same idea at work in a machine that has nothing to do with drones. The particular comes first, and the general is reached from it.

The whole book has the same shape. Part~I builds the PID controller on a drone that can only move up and down, one term and one practical problem at a time, and arrives at a full quadrotor only in its last lessons. Part~II steps up to the general language of states, models and optimisation, once the reader has seen why PID alone is not enough. The later parts of the course, on analysis and on aerospace guidance and control, take the same ideas further still. Theory always arrives when practice has asked for it.

\section*{Who it is for}
The book assumes what a bachelor's degree in engineering, physics, mathematics or computer science usually provides: calculus of one and several variables, the basics of ordinary differential equations (linear equations with constant coefficients, above all), linear algebra (matrices, eigenvalues), and a little probability. Everything beyond that — Laplace transforms, the phase margin, the Riccati equation, the Kalman filter, quadratic programming, quaternions — is introduced where it is first needed, and summarised in \hyperref[app:math]{the mathematical toolbox} at the end.

\section*{How it is organised}
The book follows the lessons of the simulator, one lesson per chapter, with the same numbers and titles you see in the application's lesson panel.
\begin{description}[leftmargin=0pt, font=\sffamily\small\bfseries]
  \item[Part I · PID] (Lessons I.1–I.14) builds the PID controller term by term on a drone that can only move up and down, meets the practical problems of a real digital controller — saturation, windup, derivative kick, noise, sampling, delay, actuator lag — and ends with the cascade of loops that flies a full quadrotor.
  \item[Part II · Beyond PID] (Lessons II.1–II.21) is a second semester in five chapters: optimal state feedback and estimation (LQR, Kalman filter), disturbance observers and incremental control (ADRC, INDI), geometry and trajectories (rotations, differential flatness, minimum-snap planning), optimisation in the loop (MPC, MPPI, control barrier functions), and adaptation and learning (L1 adaptive control, learned models, neural-network policies) — ending with a comparison of all of them on equal terms.
\end{description}
These two parts are the first volume of a longer course. Part~III, \emph{Why it works}, turns from design to analysis: frequency response, stability margins, Lyapunov's method, robustness, estimation with real sensors, and verification by Monte Carlo. Part~IV, \emph{Aerospace GNC}, takes the same ideas to a rocket, an aircraft and a satellite and adds optimal control from first principles, gain scheduling and guidance. Part~V is about engineering practice. Their lessons appear in the simulator first and in later volumes of this book; the \hyperref[ch:colophon]{colophon} says where to find the plan.

Chapter~0, \hyperref[ch:intro]{\emph{Why control?}}, tells how this book began — with a small helicopter on Mars — and then introduces the field itself: what control theory is, how it is organised, where it is used and where it came from. The \hyperref[ch:prelude]{Prelude} then introduces the simulator and its physics. Read both first, or come back to the Prelude when a lesson refers to it.

\section*{The anatomy of a lesson}
Every chapter has the same parts, set in the same way throughout the book:

\begin{keyidea}[the one thing to take away]
Boxes like this state the central idea of a section in a few sentences.
\end{keyidea}
\begin{definition}{a precise term}
Definitions fix the vocabulary.
\end{definition}
\begin{inapp}[where the idea lives in the application]
These dark boxes, styled like the application itself, say where to find an idea on the screen — \ui{Group\then Parameter} — which parameter (\param{control.alt.kp}) and which source file (\src{src/control/pid.ts}) implement it, and which keys to press (\key{R}). File names are links to the repository.
\end{inapp}
\begin{trythis}
Experiments to run in the simulator, usually with a question to answer \emph{before} running them.
\end{trythis}
\begin{watchout}
Common mistakes and the limits of a result.
\end{watchout}
\begin{example}[a problem solved in full]
Worked examples take a concrete question with the drone's real numbers and solve it step by step, with nothing left out. They are the backbone of each lesson: read them with a pencil.
\end{example}
\begin{lookahead}[the thread into the rest of the course]
A lesson closes its theory with a section \emph{The engineer's view}, which says the result again in the language used at work — state, poles, energy — and ends with a subsection \emph{In formal terms}, where the definitions and theorems are stated exactly, with their hypotheses and a proof or its source. Then comes a box like this one, which names the later lessons where the idea returns.
\end{lookahead}
Two more things recur. Wherever a formula predicts a number, a figure or a table puts the prediction next to what the simulator measured, and the text explains any difference. And each chapter ends with a page \emph{Out in the World} — the same idea at work somewhere else, with a photograph — a short summary, and exercises marked by difficulty ({\tiny$\bullet$} to {\tiny$\bullet\bullet\bullet$}). Some exercises carry a label in the margin: \exkindinline{derive} asks for a derivation with pencil and paper, \exkindinline{fly} for an experiment in the simulator, and \exkindinline{code} for reading or changing the simulator's source. The margin of a lesson's first page holds its set-up: what is switched on, and what happens when. Complete solutions of the exercises will be published in a companion volume.

\section*{One colour, one meaning}
The simulator gives every quantity its own colour, in its charts, its 3D view and its parameter panel. This book uses the same colours, so that a figure here and a chart on the screen read the same way:

\begin{center}\small\sffamily
\begin{tabular}{@{}lll@{\hspace{2.5em}}lll@{}}
  \textcolor{sigsp}{\rule[0.5ex]{1.2em}{1.2pt}} & setpoint $r$ & dashed grey &
  \textcolor{sigP}{\rule[0.3ex]{1.2em}{2pt}} & \Pt term & orange\\
  \textcolor{sigmeas}{\rule[0.3ex]{1.2em}{2pt}} & measurement $y$ & blue &
  \textcolor{sigI}{\rule[0.3ex]{1.2em}{2pt}} & \It term & green\\
  \textcolor{sigerr}{\rule[0.3ex]{1.2em}{2pt}} & error $e$ & red &
  \textcolor{sigD}{\rule[0.3ex]{1.2em}{2pt}} & \Dt term & violet\\
  \textcolor{sigwind}{\rule[0.3ex]{1.2em}{2pt}} & wind & green &
  \textcolor{sigFF}{\rule[0.3ex]{1.2em}{2pt}} & \FFt (feedforward) & amber\\
\end{tabular}
\end{center}
Red shading in a time plot means that an actuator is saturated, as in the application.

\section*{Where the numbers come from}
All time series in this book were produced by the simulator itself. A script, \src{scripts/book-data.ts}, sets up each lesson exactly as the application does, flies it headlessly with the same code, and records the telemetry; the figures are drawn from those recordings. When a number in the text says \enquote{the simulator shows \SI{0.45}{m}}, you can reproduce it. Where theory and simulation disagree — and they sometimes do, instructively — the book says so and explains why.

\section*{Running the simulator}
Anemostatos runs locally in a web browser. With Node.js installed:
\begin{center}
\texttt{git clone https://github.com/wojciech-stelmaszewski/anemostatos}\\
\texttt{cd anemostatos \&\& make dev}
\end{center}
The lesson panel on the right of the screen has two tabs, \ui{Part I · PID} and \ui{Part II · Beyond PID}; select a lesson and press \ui{Start}. Everything else — every gain, every physical parameter, the wind, the sensors — stays unlocked: the lessons are starting points, not rails.
